\section{因式分解的应用}

\Needspace{18\baselineskip}

\subsection{条件求值}

\topicline{题目只求某个整体时，先看已知条件能否直接确定它。}

\Needspace{12\baselineskip}

\begin{problem}{例题 15 （恒等式中的参数）}

设 $y=ax$，若代数式 $(x+y)(x-2y)+3y(x+y)$ 可化简为 $x^2$，求 $a$。这里“化简为”表示关于 $x$ 的恒等式。

\end{problem}

\Needspace{6\baselineskip}

\begin{solution}{例题 15 · 讲解}

\think 两项都有 $x+y$，先提出它，另一个因式也恰好是 $x+y$。代入 $y=ax$ 后，比较 $x^2$ 的系数即可，不必分别求 $x,y$。

\solve \begin{align*}
(x+y)(x-2y)+3y(x+y)&=(x+y)(x+y)\\
&=(1+a)^2x^2.
\end{align*}
由恒等式，比较系数得 $(1+a)^2=1$，即 $1+a=\pm1$，所以 $a=0$ 或 $a=-2$。
代入 $a=0$ 或 $a=-2$，原式的系数 $(1+a)^2$ 都为 $1$。

\insight{恒等式可比较系数；解平方等式时要保留两种符号。}

\end{solution}

\Needspace{12\baselineskip}

\begin{problem}{例题 16 （由两个条件求和）}

实数 $x,y$ 满足 $x^2+xy+y=14$、$y^2+xy+x=28$，求 $x+y$。

\end{problem}

\Needspace{6\baselineskip}

\begin{solution}{例题 16 · 讲解}

\think 题目只问 $x+y$，两式相加就能凑出它的平方。求出的两个和值都要保留，再用两式之差分别求出一组 $x,y$，确认原条件确实允许这些和值。

\solve 令 $t=x+y$，两式相加得
\[t^2+t=42,\qquad (t-6)(t+7)=0.\]
故 $t=6$ 或 $t=-7$。两式相减得 $(x-y)(t-1)=-14$。
当 $t=6$ 时，$x-y=-\frac{14}{5}$，得 $(x,y)=(\frac85,\frac{22}5)$；
当 $t=-7$ 时，$x-y=\frac74$，得 $(x,y)=(-\frac{21}8,-\frac{35}8)$。
将两组数分别代回原条件均成立，所以 $x+y=6$ 或 $-7$。

\insight{求出整体的候选值后，要确认它们都能满足原条件。}

\end{solution}

\methodcard{先看条件与目标式中有没有相同的整体。条件相加、相减或移项后，有时能直接求出这个整体，无须分别求每个字母。}{换元后保留原有范围；题目未说字母为正时，要保留负值的可能。}

\begin{problem}{练习 17}

实数 $a,b$ 满足 $a^4+b^4=a^2-2a^2b^2+b^2+6$，求 $a^2+b^2$。

\end{problem}

\begin{problem}{练习 18}

两个不等实数 $m,n$ 满足 $m^2-2m=a$、$n^2-2n=a$、$m^2+n^2=5$，求 $a$。

\end{problem}

\Needspace{18\baselineskip}

\subsection{简便计算}

\topicline{数字较大或因数较多时，先看各项的关系，再决定怎样分解。}

\Needspace{12\baselineskip}

\begin{problem}{例题 17 （大数计算）}

计算 $\displaystyle\frac{2004^3-2\times2004^2-2002}{2004^3+2004^2-2005}$。

\end{problem}

\Needspace{6\baselineskip}

\begin{solution}{例题 17 · 讲解}

\think 这些数都接近 $2004$。令 $t=2004$，则 $2002=t-2$、$2005=t+1$，分子、分母都能写成四项式，再用分组法分解。

\solve 令 $t=2004$，则
\begin{align*}
\text{分子}&=t^3-2t^2-t+2=(t-2)(t^2-1),\\
\text{分母}&=t^3+t^2-t-1=(t+1)(t^2-1).
\end{align*}
因 $t=2004$，上述分母和被约去的因式都不为零，故
\[\text{原式}=\frac{t-2}{t+1}=\frac{2002}{2005}.\]

\insight{大数之间有简单的和差关系时，先用字母表示，往往更容易计算。}

\end{solution}

\Needspace{12\baselineskip}

\begin{problem}{例题 18 （连乘计算）}

计算
\[
\frac{(2\times5+2)(4\times7+2)\cdots(1994\times1997+2)}
{(1\times4+2)(3\times6+2)\cdots(1993\times1996+2)}.
\]

\end{problem}

\Needspace{6\baselineskip}

\begin{solution}{例题 18 · 讲解}

\think 先看第一项：$2\times5+2=3\times4$，$1\times4+2=2\times3$，相除得到 $2$。第二组会得到 $3/2$ 吗？用第 $k$ 组的式子验证，就能看清后面的约分规律。

\solve 分子、分母各有 $997$ 项。对 $k=1,2,\ldots,997$，有
\begin{align*}
2k(2k+3)+2&=2(2k+1)(k+1),\\
(2k-1)(2k+2)+2&=2k(2k+1).
\end{align*}
所以
\[\text{原式}=\frac21\times\frac32\times\frac43\times\cdots\times\frac{998}{997}=998.\]

\insight{先分解一组对应项，再写一般项，最后核对首项、末项和项数。}

\end{solution}

\methodcard{先用字母写出数字间的关系。连乘题可从一组对应项入手，分解后再看哪些因式重复出现。}{只有分子、分母中的相同非零因式才能约去，和式中的项不能直接约。}

\begin{problem}{练习 19}

计算
\[
\frac{(7^4+64)(15^4+64)(23^4+64)(31^4+64)(39^4+64)}
{(3^4+64)(11^4+64)(19^4+64)(27^4+64)(35^4+64)}.
\]

\end{problem}

\Needspace{18\baselineskip}

\subsection{平方与非负性}

\topicline{配成平方后，除了可以分解，还可以利用平方非负判断取值。}

\Needspace{12\baselineskip}

\begin{problem}{例题 19 （平方和为零）}

已知实数 $m,n$，当 $x=m$ 时，多项式 $x^2+2x+n^2$ 的值为 $-1$；求 $x=-m$ 时该多项式的值。

\end{problem}

\Needspace{6\baselineskip}

\begin{solution}{例题 19 · 讲解}

\think 把 $x=m$ 代入，得 $m^2+2m+n^2=-1$。移项后能配成两个平方的和；和为零时，两个平方都只能为零。

\solve 由已知条件得
\[(m+1)^2+n^2=0.\]
两项均非负，故 $m+1=0$、$n=0$，即 $m=-1$。
当 $x=-m=1$ 时，原多项式的值为 $1^2+2\times1+0=3$。

\insight{实数的平方非负；几个平方之和为零时，每个平方都等于零。}

\end{solution}

\Needspace{12\baselineskip}

\begin{problem}{例题 20 （和一定，求积的最大值）}

观察 $1\times49,2\times48,\ldots,25\times25,\ldots,49\times1$。
\begin{enumerate}[label=(\arabic*)]
\item 写出两数乘积的最大值，并写出两个因数 $a,b$ 的关系。
\item 类比 $1\times59,2\times58,\ldots,59\times1$，求乘积最大值并证明。
\end{enumerate}

\end{problem}

\Needspace{6\baselineskip}

\begin{solution}{例题 20 · 讲解}

\think 第一组两因数的和总是 $50$，乘积在两数接近时较大。怎样证明？由 $4ab=(a+b)^2-(a-b)^2$，和固定后，差的平方越小，乘积越大。也可以由恒等式 $(a-b)^2=(a+b)^2-4ab\ge0$，得到 $(a+b)^2\ge4ab$。

\solve 第一组有 $a+b=50$，所以
\[ab=\frac{2500-(a-b)^2}{4}\le625.\]
当 $a=b=25$ 时等号成立，故最大值为 $625$。
第二组记两因数为 $m,n$，则 $m+n=60$，于是
\[mn=\frac{3600-(m-n)^2}{4}\le900.\]
当 $m=n=30$ 时等号成立，故最大值为 $900$。

再想一步：若两个正整数的和为 $51$，它们不可能相等，差的绝对值最小是多少？

\insight{证明最大值时，既要给出上界，也要说明等号能够取到。}

\end{solution}

\methodcard{可从交叉项决定怎样配方。配成平方后，要接着说明它对取值有什么限制：平方和为零、平方非负，或某个平方有上界。}{和为零不能直接推出每项为零。求最值时还要检查等号条件。}

\begin{problem}{练习 20}

实数 $x,y$ 满足 $x^2+12xy+52y^2-8y+1=0$，求 $x^2-y^2$。

\end{problem}

\begin{problem}{练习 21}

已知实数 $m\ne n$，$a=2m^2-mn$、$b=mn-2n^2$、$c=m^2-n^2$。用“$<$”表示 $a,b,c$ 的大小关系。

\end{problem}

\begin{problem}{练习 22}

有多少个整数解 $(x,y)$ 满足 $x^2+y^2\le2x+2y$？

\end{problem}

\Needspace{18\baselineskip}

\subsection{整数方程}

\topicline{把方程配成两个整数的乘积，常能将求解变为列举有限个因数对。}

\Needspace{12\baselineskip}

\begin{problem}{例题 21 （整数解的列举）}

求方程 $xy-2x-2y+7=0$ 的全部整数解，其中 $x\le y$。

\end{problem}

\Needspace{6\baselineskip}

\begin{solution}{例题 21 · 讲解}

\think 先把 $xy-2x-2y$ 与 $(x-2)(y-2)$ 比较，后者只多一个常数 $4$。补上这项，方程就成了两个整数的乘积等于 $-3$，可列出因数对。

\solve 原方程等价于
\[xy-2x-2y+4=-3,\qquad (x-2)(y-2)=-3.\]
整数因数对在 $x-2\le y-2$ 的限制下只有 $(-3,1)$、$(-1,3)$。
因此 $(x,y)=(-1,3)$ 或 $(1,5)$。

\insight{配成整数乘积后，先列因数对，再用大小条件筛选。}

\end{solution}

\Needspace{6\baselineskip}

\methodcard{含有 $xy$ 及 $x,y$ 一次项的方程，可尝试添一个常数配成乘积。例如 $xy+x+y=M$ 可写成 $(x+1)(y+1)=M+1$（回看练习 7、8 后的总结）。若 $x,y$ 为正整数，新因数都至少为 $2$；这些范围要一并用来筛选。}{整数可以为负。}

\begin{problem}{练习 23}

整数 $a,b$ 满足 $6ab=9a-10b+303$，求 $a+b$。

\end{problem}

\begin{problem}{练习 24}

从 $1,2,\ldots,37$ 中选出两个不同的数，使其余 $35$ 个数之和等于这两个数之积。求所选两数之差的绝对值。

\end{problem}

\begin{problem}{练习 25}

两个互不相等的正整数 $a,b$ 满足 $a^2-b^2=2018-2a$，求 $a+b$。

\end{problem}

\Needspace{18\baselineskip}

\subsection{数的性质：整除、质数与完全平方}

\topicline{先分解，再判断整除性；遇到质数或完全平方数条件，还要利用因数的范围。}

\Needspace{12\baselineskip}

\begin{problem}{例题 22 （用整除性筛选）}

将某个正整数 $n$ 代入 $n^5-n$，四个结果中仅一个正确。它是哪一个？A．$7770$；B．$7775$；C．$7776$；D．$7779$。

\end{problem}

\Needspace{6\baselineskip}

\begin{solution}{例题 22 · 讲解}

\think 不必直接算五次方。先分解 $n^5-n$，看它一定能被哪些数整除，再据此排除选项。选出一个结果后，还需找到相应的 $n$。

\solve \[n^5-n=n(n^4-1)=n(n-1)(n+1)(n^2+1).\]
连续三个整数中有一个偶数、一个 $\boxed{3\text{ 的倍数}}$，因此该式必被 $6$ 整除。
先排除奇数选项 B、D。A、C 都是 $6$ 的倍数，还不能仅凭这一点确定答案。
再看因数 $5$：$n$ 除以 $5$ 的余数为 $0,1,2,3,4$ 时，分别有 $n,n-1,n^2+1,n^2+1,n+1$ 被 $5$ 整除。
故该式必被 $30$ 整除，四个选项中只有 A 符合。
取 $n=6$，$6^5-6=7770$，所以 A 确实可以取到。

\insight{整除性可以排除选项；要确认某个值可取到，还应给出对应的整数。}

\end{solution}

\Needspace{12\baselineskip}

\begin{problem}{例题 23 （完全平方数条件）}

自然数 $n$ 满足 $n^2+15n+26$ 是完全平方数，求 $n$。

\end{problem}

\Needspace{6\baselineskip}

\begin{solution}{例题 23 · 讲解}

\think 设原式为 $t^2$。因为一次项系数 $15$ 为奇数，先乘 $4$ 再配方，能避开分数，得到两个整数平方之差为 $121$，再列因数对。

\solve 设 $t$ 是非负整数，且 $t^2=n^2+15n+26$。则
\[(2n+15)^2-(2t)^2=121,\]
即 $(2n+15-2t)(2n+15+2t)=121$。
两因数的积为正，和为 $4n+30>0$，所以均为正整数。又前者不大于后者，因此只有 $(1,121)$、$(11,11)$ 两对。
第一对给出 $2n+15=61$，得 $n=23$、$t=30$；第二对给出 $n=-2$，不符合自然数条件。
故 $n=23$，此时原式为 $900=30^2$。

\insight{一次项系数为奇数时，可先乘 $4$，把配方写成整数形式。}

\end{solution}

\Needspace{12\baselineskip}

\begin{problem}{例题 24 （最大公因数）}

对所有大于 $2$ 的正整数 $n$，求这些整数 $n^5-5n^3+4n$ 的最大公因数。

\end{problem}

\Needspace{6\baselineskip}

\begin{solution}{例题 24 · 讲解}

\think 先提出 $n$，余下部分按 $n^2$ 分解，会得到五个连续整数的积。先证明每个值都是 $120$ 的倍数，再找一个恰好等于 $120$ 的值，才能确定最大公因数。

\solve \[n^5-5n^3+4n=n(n^2-1)(n^2-4)=(n-2)(n-1)n(n+1)(n+2).\]
五个连续整数中至少有一个是 $5$ 的倍数、一个是 $3$ 的倍数。
其乘积也被 $8$ 整除：若 $n$ 为奇数，$n-1,n+1$ 是相邻偶数，其中一个被 $4$ 整除；若 $n$ 为偶数，$n-2,n,n+2$ 中有一个被 $4$ 整除，另有一个偶数提供因数 $2$。
因为 $3,5,8$ 两两互质，乘积被 $120$ 整除。
取 $n=3$，原式恰为 $1\times2\times3\times4\times5=120$，故最大公因数为 $120$。

\insight{求最大公因数，要证明它是公因数，还要说明公因数不能更大。}

\end{solution}

\Needspace{12\baselineskip}

\par\noindent\begin{minipage}{\linewidth}

\begin{problem}{例题 25 （四个连续整数的积加一）}

设 $n$ 为正整数，证明 $(n+2020)(n+2021)(n+2022)(n+2023)+1$ 是完全平方数。

\end{problem}

\Needspace{6\baselineskip}

\begin{solution}[breakable=false]{例题 25 · 讲解}

\think 令 $N=n+2020$，四个因数就成为连续整数。先把首尾一对、中间一对相乘，得到 $N^2+3N$ 与 $N^2+3N+2$。这与例题 5 的两个因子相差 $2$ 的情形相同。

\solve 令 $N=n+2020$，则

\begin{align*}
\text{原式}&=N(N+1)(N+2)(N+3)+1\\
&=[N(N+3)][(N+1)(N+2)]+1\\
&=(N^2+3N)(N^2+3N+2)+1\\
&=(N^2+3N+1)^2.
\end{align*}

因为 $N^2+3N+1$ 是整数，所以原式是完全平方数。更一般地，任意整数 $m$ 都满足

\[
m(m+1)(m+2)(m+3)+1=(m^2+3m+1)^2.
\]

\insight{四个连续整数对应的式子可以两两配对，使两组乘积含有相同的部分；两组乘积相差 $2$，加一后可配成平方。}

\end{solution}

\end{minipage}\par\medskip

\methodcard{整除题可先分解，再找所需的整数因数。质数条件常限制某个正因数为 $1$；完全平方条件可设平方后配成平方差，也可直接把原式配成平方。求最大公因数还应找一个取值来限制它不能更大。}{分别能被 $a,b$ 整除，只有在 $a,b$ 互质时，才能直接推出能被 $ab$ 整除。}

\begin{problem}{练习 26}

已知 $n$ 为正整数，且 $n^4-16n^2+100$ 是质数，求 $n$。

\end{problem}

\begin{problem}{练习 27}

证明 $2018^2+2018^2\times2019^2+2019^2$ 是完全平方数。

\end{problem}

\Needspace{30\baselineskip}

\subsection{几何中的数量关系}

\topicline{把图中的长度与面积写成代数式，再用已学的恒等变形处理。}

\Needspace{12\baselineskip}

\begin{problem}{例题 26 （图形中的面积关系）}

如图，分别以正数 $a,b,m,n$ 为边长作正方形，$m>n$，且 $am-bn=2$、$an+bm=4$。
图①是边长为 $a,b$ 的两个正方形；图②中 $AB=m$、$DC=n$、$BP=m$、$PC=n$，$AB,DC$ 均垂直于 $BC$，阴影为 $\triangle APD$。
\begin{enumerate}[label=(\arabic*)]
\item 若 $a=3,b=4$，求图①涂色面积。
\item 若图①涂色面积为 $3$，四边形 $ABCD$ 面积为 $5$，求图②阴影面积。
\end{enumerate}
\factorfigure

\end{problem}

\Needspace{6\baselineskip}

\begin{solution}{例题 26 · 讲解}

\think 两式平方相加时，含 $abmn$ 的项会抵消，剩下 $(a^2+b^2)(m^2+n^2)$。图①给出 $a^2+b^2$，梯形面积给出 $(m+n)^2$；阴影面积可由梯形减去两个直角三角形求得。

\solve 第 (1) 问：涂色面积为 $a^2+b^2=3^2+4^2=25$。
第 (2) 问：展开并分组得
\begin{align*}
(am-bn)^2+(an+bm)^2&=a^2(m^2+n^2)+b^2(m^2+n^2)\\
&=(a^2+b^2)(m^2+n^2)=20.
\end{align*}
由 $a^2+b^2=3$ 得 $m^2+n^2=\frac{20}{3}$。
梯形两条平行边是 $AB,DC$，它们之间的距离为 $BC=m+n$，故
\[S_{ABCD}=\frac{(m+n)^2}{2}=5.\]
于是 $2mn=(m+n)^2-(m^2+n^2)=10-\frac{20}{3}=\frac{10}{3}$。
又 $S_{ABP}=\frac12m^2$、$S_{PCD}=\frac12n^2$，因此
\[S_{APD}=S_{ABCD}-S_{ABP}-S_{PCD}=5-\frac12\cdot\frac{20}{3}=\frac53.\]

\insight{先写出图中各部分的面积，再看已知的代数关系能求出哪些量。}

\end{solution}

\methodcard{面积关系可以给出平方和、和的平方或乘积。若三边关系中同时有平方项和两两乘积，可尝试配成几个差的平方。}{先确认面积式对应的是图中哪一部分，尤其要找准梯形的高。}

\Needspace{12\baselineskip}

\subsection{选学挑战}

\topicline{这两题需要把前面的方法用于新的情境，可以先试算、记录，再寻找一般的理由。}

\begin{problem}{挑战 1 \starred （黑板上的操作，选学）}

黑板上有 $1,\frac12,\ldots,\frac1{100}$ 共 $100$ 个数。每次删去两个数 $a,b$，写入 $a+b+ab$。经过 $99$ 次操作后，剩下的数是多少？

\end{problem}

\methodcard{操作题可从一次操作前后的等式找不变量。同值条件可先作差，再看所得关系能否用于目标式。}{试算有助于猜想；要说明每次操作或所有取值都成立，还需给出一般理由。}

\begin{problem}{挑战 2 \starred （四次式的同值条件，选学）}

实数 $a\ne b$、正实数 $c$ 满足 $a^4-2019a=b^4-2019b=c$。证明 $-\sqrt c<ab<0$。

\end{problem}

\Needspace{12\baselineskip}
\block{做完本章，回头想一想}
\begin{enumerate}[label=(\arabic*)]
\item 比较例题 16 与练习 17：为什么前者的负值要保留，后者的负值要舍去？依据来自哪里？
\item 比较例题 22 与例题 24：两个五次式分别怎样分解？连续整数的个数相同吗？
\item 从练习 23—25 中选一道，说明每个条件在筛选因数对时用在了哪一步。再合上解析，独立重做一道曾经需要提示的题。
\end{enumerate}

换序自测可按练习 $18,26,21,24,27$ 的顺序进行。每题先判断：题目中的哪条信息提示你采用这种方法？看过解析才做出的题，隔一天再试。
